\documentclass[aspectratio=1610]{beamer}

\title{Formal Guided Test Sequence Optimization}
\date{Processor Verification - 29 May 2026}
\author{Lorenzo Pattaro Zonta}

% Theme options — combine any of these in a single \usetheme call:
   \usetheme[
    pattern=squares,
%     color=nudarkblue,
%     font=poppins,
    numbering=number,
    sectionpage=none,
%     agendastyle=pattern
    breadcrumb=on
   ]{wildcat}
%
% pattern: facet (default), facetrand, triangle, triangleiso, trianglemesh,
%          squares, circles, hexraise, hexrand, network, quilt, plain
% color: any named color (nu* colors are all available after theme loads).
%        All wcprimary shades are auto-generated. Omit to keep NU purple.
% font: noto (Noto Sans), poppins (Poppins+IBM Plex), nu (Campton+Akkurat),
%        installed (system fonts), or omit for default (Fira Sans — no files needed)
% numbering: fraction (default), number, progressbar, none
% sectionpage: on (default) or none
% agendastyle: plain (default) or pattern
% breadcrumb: on or none (default)
% \usetheme{wildcat}

% You change the titlegraphic to whatever you want, or comment it out to remove it.
\titlegraphic{\includegraphics[scale=0.25]{1upc-blanc-fons-transp.png}}

% % For fine-grained control over shades, override colors manually AFTER \usetheme.
% % For example, Yale Blue (#00356b) with hand-tuned shades:
% \definecolor{wcprimary}{RGB}{0,53,107}
% \definecolor{wcprimary140}{RGB}{0, 34, 70}
% \definecolor{wcprimary130}{RGB}{0, 40, 80}
% \definecolor{wcprimary120}{RGB}{0, 45, 91}
% \definecolor{wcprimary110}{RGB}{0, 50, 102}
% \definecolor{wcprimary40}{RGB}{153, 174, 196}
% \definecolor{wcprimary30}{RGB}{179, 194, 211}
% \definecolor{wcprimary20}{RGB}{204, 215, 225}
% \definecolor{wcprimary10}{RGB}{230, 235, 240}

% % Alert and example colors:
% \definecolor{wcalerted}{RGB}{189,83,25}
% \definecolor{wcexample}{RGB}{95,113,45}

% % If you want to change the general slide background color), 
% % you can use the following command:
%\setbeamercolor{background canvas}{bg=nupurple10!30}

% % Turn off section slides (pass sectionpage=none to \usetheme)
% \usetheme[sectionpage=none]{wildcat}

% % Use a custom background pattern 
% % (see bg-gallery/bg-template.tex for an example 
% % of how to define your own pattern):
% \input{bg-gallery/bg-template}

\begin{document}

\begin{frame}
\titlepage
\end{frame}

\begin{frame}{Introduction}
Verification process faces significant challenges:
        \begin{itemize}
        \item<2-> \# of test sequences grows exponentially with \# of test cases.
        \item<2-> Executing all possible sequences is often infeasible % due to time and resource constraints.
        \item<2-> Identification of the most effective test sequences: crucial for efficient testing.
    \end{itemize}
\end{frame}

\begin{frame}[plain]
    \tableofcontents
\end{frame}

\section{Current Verification Approaches}
\standout{Current Verification Approaches}



\begin{frame}{Current Verification Approaches}
    \textbf{Formal verification}:
    \begin{itemize}
        \item<2-> Comprehensive design exploration.
        \item<2-> Encounters state space explosion problem for large scale designs. 
    \end{itemize}
    \textbf{Simulation-based verification}:
    \begin{itemize}
        \item<3-> High scalability and applicability.
        \item<3-> Effectiveness rely on well-designed constraints.
        \item<3-> Biased constraints: miss critical scenarios.
        \item<3-> Strict constraints: increase debugging complexity.
    \end{itemize}
\end{frame}

\begin{frame}{Formal Guided Test Sequence Optimization}
    \begin{itemize}
% ● Enhance coverage & verification efficiency
% ● Leverages simulation scalability and formal verification precision to target critical
% conditions, reducing simulation overhead and improve coverage
% ● Challenges:
% ○ Resolve false alarms from formal simulation discrepancies
% ○ Address BBOX constraints
% ● Integrate simulation and formal verification
% ● Introduce abstract models to overcome limitations of BBOX
        \item<1-> Enhance coverage and verification efficiency.
        \item<2-> Leverage simulation scalability and formal verification precision to target critical conditions %, reducing simulation overhead and improving coverage.
        \item<3-> Challenges:
        \begin{itemize}
            \item<4-> Resolve false alarms from formal simulation discrepancies.
            \item<4-> Address BBOX constraints.
        \end{itemize}
        \item<5-> Integrate simulation and formal verification.
        \item<6-> Introduce abstract models to overcome limitations of BBOX.
    \end{itemize}
\end{frame}

% Both these approaches lack methodologies to ensure comprehensive coverage of critical design aspects, as in complex control flows, where formal guidance is essential

\section{Related Works}

\standout{Related Works}

\begin{frame}{Related Works}
RISC-V Torture Test:
    \begin{itemize}
        \item<2-> Generate random instruction sequences to achieve faster coverage closure
        \item<2-> Lacks structured verification guidance: limited effectiveness in extreme corner cases
    \end{itemize}
RISCV-DV:
    \begin{itemize}
        \item<3-> Support CRV (constraint random verification) and RISC V ISA features
        \item<3-> Reliance on randomness makes it difficult to cover boundary conditions
    \end{itemize}
Fuzzing tools
\begin{itemize}
    \item<4->Cannot guarantee coverage of all possible scenarios.
\end{itemize}

\end{frame}

\begin{frame}{Problems \& Issues}
\begin{itemize}
    \item False alarms due to environmental discrepancies.
    \item BBOX issues by excessive complexity or large storage modules.
\end{itemize}
\end{frame}

\section{FGTSO Framework}
\standout{Formal Guided Test Sequence Optimization Framework}

\begin{frame}{Overall architecture}
    \begin{figure}
        \centering
        \includegraphics[width=0.8\textwidth]{FGTSO_arch.png}
        \caption{FGTSO Framework Architecture}
        \label{fig:architecture}
    \end{figure}
\end{frame}

\begin{frame}{Overview}
    \begin{enumerate}
        \item<1-> Coverage Hole Selection
        \item<2-> Environment Consistency Check
        \item<3-> Formal Cover Property Generation (\textit{assume} properties)
        \item<4-> Targeted Instruction Sequence Generation
    \end{enumerate}
\end{frame}

\begin{frame}{FGTSO Workflow}   
    \begin{figure}
        \centering
        \includegraphics[width=0.4\textwidth]{workflow.png}
        \caption{FGTSO Workflow}
        \label{fig:workflow}
    \end{figure}

\end{frame}


\begin{frame}{FGTSO Workflow - Step 1}
    \begin{figure}
        \centering
        \includegraphics[width=0.8\textwidth]{step1.png}
        \caption{Step 1: Coverage Hole Selection}
        \label{fig:step1}
    \end{figure}

\end{frame}

\begin{frame}{FGTSO Workflow - Step 1}
    % The strategy prioritizes addressing the most impactful coverage holes, by focusing on
% the ones that affect verification efficiency. This accelerates the verification process.
% The paper employs the MaxUncovd strategy that prioritizes selecting modules with
% the lowest coverage, assuming that addressing the most uncovered areas yields the
% greatest verification impact. MaxUncovd enhances the overall efficiency and ensures
% comprehensive design verification.
\begin{itemize}
    \item<1-> Priority: addressing \textbf{most impactful coverage holes} % by focusing on the ones that affect verification efficiency, accelerating the verification process.
    \item<2-> MaxUncovd strategy: prioritizes selecting modules with lowest coverage, assuming that addressing the most uncovered areas yields the greatest verification impact. % MaxUncovd enhances the overall efficiency and ensures comprehensive design verification.
\end{itemize}

\end{frame}

\begin{frame}{FGTSO Workflow - Step 2}
    \begin{figure}
        \centering
        \includegraphics[width=0.8\textwidth]{step2.png}
        \caption{Step 2: Environment Consistency Check}
        \label{fig:step2}
    \end{figure}
\end{frame}

\begin{frame}{FGTSO Workflow - Step 2}
    % Misalignment detection between formal verification env with the simulation env is
% performed to determine the consistency between the two envs. Formal verification
% generates instruction sequences targeting specific coverage holes, executed in the
% simulation. If the results from formal verification and simulation are inconsistent (fail
% to cover intended coverage holes) it indicates a misalignment.
    \begin{itemize}
        \item<1-> Misalignment detection is performed to determine consistency.
        \item<2-> Formal verification generates instruction sequences targeting specific coverage holes, executed in the simulation.
        \item<3-> Results from formal verification and simulation are inconsistent (fail to cover intended coverage holes): \textbf{misalignment}.
    \end{itemize}
\end{frame}

\begin{frame}{FGTSO Workflow Misalignment - PIs}
    Three are the main causes of misalignment: PIs, BBOX and Insufficient RISCV-DV constraints.
\begin{itemize}
%     Primary inputs (any signal that originates outside the boundary of the
% specific hardware module you are testing and flows into it, e.g. reset, clock,
% interrupt request, rvalid signal [read valid] for AXI protocol). By writing
% "assumptions," verification engineers tie these wildcards down to reality so
% the formal tool only explores states that can actually happen.). PIs are treated
% as free nets, allowing them to assume arbitrary values without constraints.
% This can lead to incorrect coverage analysis, as generated cover traces may
% fail to address coverage holes in simulation. Furthermore, unconstrained PIs
% can result in misleading coverage estimation, especially when test scenarios
% correspond to unreachable states by design (es AXI protocol). Formal
% verification may fail to generate meaningful instruction patterns, leaving
% coverage holes unresolved. Formal assumptions must be introduced to align
% formal verification with realistic system behavior.
    \item<1-> Primary Inputs (PIs) (reset, clock, interrupt request, rvalid signal AXI protocol). 
    \begin{itemize}
        \item<2-> Assume arbitrary values without constraints.
        \item<2-> Generated cover traces may fail to address coverage holes.
        \item<2-> Unconstrained PIs can result in misleading coverage estimation (unreachable states by design)
        \item<2-> Formal assumptions \textbf{must be introduced} to align formal verification with system behavior.
    \end{itemize}
\end{itemize}

\end{frame}

\begin{frame}{FGTSO Workflow Misalignment - BBOX Constraints}
\begin{itemize}
    \item<2-> Modules with high storage requirements or complex computational logic: treated as BBOX.
    \item<2-> BBOX may become free nets.
    \item<2-> \textbf{Solution}: Abstraction models preserve critical behavior while ensuring interaction with the system remains verifiable.
    \item<2-> \textbf{High Storage Example}: Cache abstraction using SVA array to track indices. Predefined tag replaces full tag tracking
    \item<2-> \textbf{Computation Example}: Symbolic computation approach using SVA.
\end{itemize}

\end{frame}

\begin{frame}{Cache Abstraction Model}
    \begin{figure}
        \centering
        \includegraphics[width=0.8\textwidth]{cache_abstraction.png}
        \caption{Cache Abstraction Model}
        \label{fig:cache-abstraction}
    \end{figure}
\end{frame}

\begin{frame}{Cache Abstraction Model Codification}
    \begin{figure}
        \centering
        \includegraphics[width=1\textwidth]{cache abstraction code.png}
        \caption{SVA Implementation of Cache Abstraction}
        \label{fig:cache-code}
    \end{figure}
\end{frame}


\begin{frame}{FGTSO Workflow Misalignment - RISCV-DV Constraints}
\begin{itemize}
    \item<1-> \textbf{Problem}: Formal verification may produce sequences that RISCV-DV constraints cannot handle.
    \item<2-> \textbf{Example}: Multi-level page table mechanism includes 1 GB, 2 MB, and 4 KB page tables, but original constraints only allow 1 GB page table access.
    \item<3-> \textbf{Solution}: Modify constraints to apply them to addresses pointed to by 1 GB and 2 MB page table entries.
    \item<4-> \textbf{Implementation}: Set X, W, and R bits to 0 to indicate that entries should point to the next level of the page table.
    \item<5-> \textbf{Result}: Coverage holes for multi-level page table accesses can now be filled.
\end{itemize}

\end{frame}

\begin{frame}{Step 3 \& 4: Formal Cover Property Generation}
\begin{itemize}
    \item<2-> \textbf{Overview}: Use formal cover properties to obtain cover traces and generate instruction sequences addressing target coverage holes.
    \item<2-> \textbf{Property Definition}: Step automatically extracts coverage hole expression and applies it to predefined cover templates.
    \item<2-> \textbf{Formal Execution}: Formal tool executes the cover property and derives the input patterns required to satisfy the condition.
    \item<2-> \textbf{Cover Trace Generation}: Generates corresponding cover trace from the formal verification results.
    \item<2-> \textbf{Sequence Extraction}: Extract the instruction sequence from the trace to effectively cover the coverage hole.
\end{itemize}

\end{frame}

\begin{frame}{Coverage Holes Analysis}
    \begin{figure}
        \centering
        \includegraphics[width=0.85\textwidth]{Coverage Holes.png}
        \caption{Coverage Holes Identification}
        \label{fig:coverage-holes}
    \end{figure}
\end{frame}

\begin{frame}{Formal Cover Property Example}
    \begin{figure}
        \centering
        \includegraphics[width=0.85\textwidth]{cover_property.png}
        \caption{Cover Property Definition}
        \label{fig:cover-property}
    \end{figure}
\end{frame}

\begin{frame}{Handling Coverage Challenges}
\begin{itemize}
    \item<2-> \textbf{Unreachable Coverage Holes}:constraints are overly restrictive?
    \item<2-> \textbf{Waiving Unreachable Holes}: If appropriate \& hole is unreachable, it is waived; information still logged for debugging.
    \item<2-> \textbf{Resource Exhaustion}: State space explosion.
    \item<2-> \textbf{Decomposition Strategy}: Complex properties decomposed in smaller sub-properties
    \item<2-> \textbf{Integration}: Integrate sub-property traces into comprehensive instruction sequence
\end{itemize}

\end{frame}

\begin{frame}{Final Outcome}
    \begin{figure}
        \centering
        \includegraphics[width=0.9\textwidth]{results.png}
        \caption{FGTSO Results: Coverage Achievement}
        \label{fig:results}
    \end{figure}
\end{frame}

\section{Results}
\standout{Results}
\begin{frame}{Comparison with Other Frameworks}
    \begin{figure}
        \centering
        \includegraphics[width=0.9\textwidth]{framework comparison.png}
        \caption{FGTSO vs. Other Verification Approaches}
        \label{fig:framework-comparison}
    \end{figure}
\end{frame}

\begin{frame}{100\% Coverage}
    \begin{figure}
        \centering
        \includegraphics[width=0.9\textwidth]{FGTSO Coverage.png}
        \caption{FGTSO: Achieving Complete Coverage}
        \label{fig:fgtso-coverage}
    \end{figure}
\end{frame}

\standout{Thanks for your attention!}

\begin{frame}
    \vspace{3cm}
    \begin{center}
        \Large Questions?
    \end{center}
\vspace{2.5cm}
    \begin{pblock}{Based on the paper}
        SHEN, Yu-Shien, et al. Efficient Coverage Optimization with Formal-Guided Testcase Generation in UVM Verification. In: DVCon Europe 2025; Design and Verification Conference and Exibition. VDE, 2025. p. 43-48.
    \end{pblock}

\end{frame}

\appendix

\end{document}
